\documentclass[11pt,a4paper]{article}
\usepackage[utf8]{inputenc}
\usepackage[vietnamese]{babel}
\usepackage{amsmath,amssymb}
\usepackage[version=4]{mhchem}
\usepackage{geometry}
\geometry{top=1.5cm,bottom=1.2cm,left=1.4cm,right=1.4cm,headheight=28pt,headsep=12pt,footskip=18pt}
\usepackage{graphicx}
\usepackage{xcolor}
\usepackage{tabularx}
\usepackage{array}
\usepackage{colortbl}
\usepackage{fancyhdr}
\usepackage{enumitem}

\definecolor{hfRed}{HTML}{D32F2F}
\definecolor{hfGreen}{HTML}{8BC390}
\definecolor{mainGreen}{HTML}{2E7D32}
\definecolor{solBlue}{HTML}{1565C0}
\definecolor{tableGreen}{HTML}{E8F5E9}
\definecolor{ansRed}{HTML}{C62828}

\pagestyle{fancy}
\fancyhf{}
\renewcommand{\headrulewidth}{0.75pt}
\renewcommand{\headrule}{\hbox to\headwidth{\color{hfGreen}\leaders\hrule height \headrulewidth\hfill}}
\renewcommand{\footrulewidth}{0.75pt}
\renewcommand{\footrule}{\hbox to\headwidth{\color{hfGreen}\leaders\hrule height \footrulewidth\hfill}}

\fancyhead[L]{\small\itshape\color{hfRed}Tài liệu lưu hành nội bộ}
\fancyhead[R]{\small\color{hfRed}Thầy \textbf{TRẦN VĂN THẠNH} - 0777.470.803}
\fancyfoot[L]{\footnotesize\itshape\color{hfRed}CS1: 6/15 Nguyễn Hoàng, P. Kim Long, TP Huế \quad | \quad CS2: 24 Đặng Thái Thân, TP Huế}
\fancyfoot[R]{\small\bfseries\color{hfRed}Trang \thepage}

\setlength{\parindent}{0pt}
\setlength{\parskip}{3pt}

\newcommand{\dapan}[1]{\textcolor{ansRed}{\textbf{#1}}}
\newcommand{\loigiai}[1]{
    \par\vspace{2pt}
    \begingroup\leftskip=12pt
    \textcolor{solBlue}{\textbf{$\blacktriangleright$ Lời giải chi tiết:}}\par
    \small #1
    \par\endgroup\vspace{3pt}
}

\begin{document}


\begin{table}[h!]
\centering
\begin{tabularx}{\textwidth}{@{}p{8.5cm}X@{}}
\textbf{SỞ GD\&ĐT TP HUẾ} & \textbf{KIỂM TRA GIỮA KÌ I NĂM HỌC 2026 - 2027} \\
\textbf{TRƯỜNG THPT HAI BÀ TRƯNG} & \textbf{MÔN: HÓA HỌC 12} \\
\textbf{ĐỀ CHÍNH THỨC} & Thời gian làm bài: 45 phút \\
\textit{(Đề thi có 03 trang)} & \textbf{MÃ ĐỀ: 101} \\
\end{tabularx}
\end{table}


\noindent\textbf{Họ, tên thí sinh:} \dotfill \textbf{Lớp:} \dots\dots\dots\dots \textbf{SBD:} \dots\dots\dots\dots\\
\textit{(Cho biết nguyên tử khối: H = 1; C = 12; N = 14; O = 16; Na = 23; K = 39; Cl = 35,5; Ag = 108).}
\par\vspace{-2pt}\noindent\textcolor{gray!40}{\rule{\textwidth}{0.5pt}}\par\vspace{2pt}


\noindent\textbf{\large\color{mainGreen}PHẦN I. Câu trắc nghiệm nhiều phương án lựa chọn (3,0 điểm)}\\
\textit{Thí sinh trả lời từ Câu 1 đến Câu 12. Mỗi câu hỏi thí sinh chỉ chọn một phương án.}
\vspace{4pt}

\textbf{Câu 1.} Chất nào sau đây thuộc loại ester?\\[2pt]
\begin{tabularx}{\textwidth}{@{}XXXX@{}}\textbf{A.} \ce{CH3COOH}. & \textbf{B.} \ce{CH3COOCH3}. & \textbf{C.} \ce{C2H5OH}. & \textbf{D.} \ce{CH3CHO}.\end{tabularx}\\[2pt]
\vspace{2pt}
\textbf{Câu 2.} Thành phần chính của xà phòng thường là\\[2pt]
\begin{tabularx}{\textwidth}{@{}XX@{}}\textbf{A.} muối sodium hoặc potassium của acid béo. & \textbf{B.} muối sodium alkylsulfate. \\ \textbf{C.} muối sodium alkylbenzene sulfonate. & \textbf{D.} triester của glycerol với acid béo.\end{tabularx}\\[2pt]
\vspace{2pt}
\textbf{Câu 3.} Carbohydrate nào sau đây thuộc loại monosaccharide?\\[2pt]
\begin{tabularx}{\textwidth}{@{}XXXX@{}}\textbf{A.} Saccharose. & \textbf{B.} Tinh bột. & \textbf{C.} Glucose. & \textbf{D.} Cellulose.\end{tabularx}\\[2pt]
\vspace{2pt}
\textbf{Câu 4.} Ở dạng mạch hở, trong phân tử glucose có chứa nhóm chức nào sau đây?\\[2pt]
\begin{tabularx}{\textwidth}{@{}XX@{}}\textbf{A.} Chỉ có nhóm hydroxy (-OH). & \textbf{B.} Một nhóm ketone (-CO-) và năm nhóm -OH. \\ \textbf{C.} Một nhóm aldehyde (-CHO) và năm nhóm -OH. & \textbf{D.} Một nhóm carboxyl (-COOH) và bốn nhóm -OH.\end{tabularx}\\[2pt]
\vspace{2pt}
\textbf{Câu 5.} Cho glucose tác dụng với lượng dư thuốc thử Tollens ([\ce{Ag(NH3)2}]OH), đun nóng nhẹ. Hiện tượng quan sát được là\\[2pt]
\quad \textbf{A.} xuất hiện kết tủa màu đỏ gạch.\\[2pt]
\quad \textbf{B.} có lớp kim loại bạc sáng bóng bám vào thành ống nghiệm.\\[2pt]
\quad \textbf{C.} dung dịch chuyển sang màu xanh lam đậm.\\[2pt]
\quad \textbf{D.} có sủi bọt khí không màu thoát ra.\\[2pt]
\vspace{2pt}
\textbf{Câu 6.} Chất nào sau đây là thành phần chính của đường mía, đường phèn thường dùng trong pha chế đồ uống hàng ngày?\\[2pt]
\begin{tabularx}{\textwidth}{@{}XXXX@{}}\textbf{A.} Glucose. & \textbf{B.} Fructose. & \textbf{C.} Saccharose. & \textbf{D.} Maltose.\end{tabularx}\\[2pt]
\vspace{2pt}
\textbf{Câu 7.} Chất nào sau đây thuộc loại amine bậc một?\\[2pt]
\begin{tabularx}{\textwidth}{@{}XXXX@{}}\textbf{A.} \ce{CH3-NH-CH3}. & \textbf{B.} \ce{(CH3)3N}. & \textbf{C.} \ce{C2H5NH2}. & \textbf{D.} \ce{C6H5-NH-CH3}.\end{tabularx}\\[2pt]
\vspace{2pt}
\textbf{Câu 8.} Ở điều kiện thường, amine nào sau đây tồn tại ở thể khí?\\[2pt]
\begin{tabularx}{\textwidth}{@{}XXXX@{}}\textbf{A.} Aniline. & \textbf{B.} Methylamine. & \textbf{C.} Benzylamine. & \textbf{D.} Hexan-1-amine.\end{tabularx}\\[2pt]
\vspace{2pt}
\textbf{Câu 9.} Nhỏ vài giọt nước bromine vào ống nghiệm chứa aniline, hiện tượng quan sát được là\\[2pt]
\begin{tabularx}{\textwidth}{@{}XX@{}}\textbf{A.} xuất hiện kết tủa trắng. & \textbf{B.} dung dịch chuyển sang màu xanh lam. \\ \textbf{C.} có bọt khí bay lên. & \textbf{D.} xuất hiện kết tủa màu đỏ gạch.\end{tabularx}\\[2pt]
\vspace{2pt}
\textbf{Câu 10.} Công thức cấu tạo thu gọn của Alanine là\\[2pt]
\begin{tabularx}{\textwidth}{@{}XX@{}}\textbf{A.} \ce{H2N-CH2-COOH}. & \textbf{B.} \ce{CH3-CH(NH2)-COOH}. \\ \textbf{C.} \ce{H2N-(CH2)2-COOH}. & \textbf{D.} \ce{HOOC-CH2-CH(NH2)-COOH}.\end{tabularx}\\[2pt]
\vspace{2pt}
\textbf{Câu 11.} Số liên kết peptide có trong phân tử tripeptide Gly-Ala-Val là\\[2pt]
\begin{tabularx}{\textwidth}{@{}XXXX@{}}\textbf{A.} 3. & \textbf{B.} 2. & \textbf{C.} 1. & \textbf{D.} 4.\end{tabularx}\\[2pt]
\vspace{2pt}
\textbf{Câu 12.} Chất nào sau đây thuộc loại protein hình cầu, tan được trong nước tạo dung dịch keo?\\[2pt]
\begin{tabularx}{\textwidth}{@{}XX@{}}\textbf{A.} Keratin (trong tóc, móng). & \textbf{B.} Albumin (trong lòng trắng trứng). \\ \textbf{C.} Fibroin (trong tơ tằm). & \textbf{D.} Collagen (trong da, gân).\end{tabularx}\\[2pt]
\vspace{2pt}

\vspace{4pt}
\noindent\textbf{\large\color{mainGreen}PHẦN II. Câu trắc nghiệm đúng sai (2,0 điểm)}\\
\textit{Thí sinh trả lời từ Câu 1 đến Câu 2. Trong mỗi ý a), b), c), d) ở mỗi câu, thí sinh chọn đúng hoặc sai.}
\vspace{4pt}


\textbf{Câu 1.} Thực hiện thí nghiệm xà phòng hóa chất béo trong phòng thí nghiệm:\\
\textit{Bước 1:} Cho vào bát sứ nhỏ khoảng 1 gam mỡ lợn (chứa tripalmitin) và 2,5 mL dung dịch NaOH 40\%.\\
\textit{Bước 2:} Đun sôi nhẹ hỗn hợp, liên tục khuấy đều bằng đũa thủy tinh trong khoảng 8 - 10 phút và thỉnh thoảng thêm vài giọt nước cất.\\
\textit{Bước 3:} Rót thêm vào hỗn hợp 5 mL dung dịch NaCl bão hòa nóng, khuấy nhẹ, sau đó để nguội.\\[2pt]

\quad \textbf{a)} Tripalmitin là triester của glycerol với palmitic acid, có công thức là \ce{(C15H31COO)3C3H5}.\\[2pt]
\quad \textbf{b)} Ở điều kiện thường, tripalmitin là chất lỏng, nhẹ hơn nước, không tan trong nước nhưng tan nhiều trong các dung môi hữu cơ không phân cực.\\[2pt]
\quad \textbf{c)} Mục đích chính của việc thêm dung dịch NaCl bão hòa ở bước 3 là làm cho xà phòng (muối sodium của acid béo) tách ra khỏi dung dịch, nổi lên trên mặt dung dịch.\\[2pt]
\quad \textbf{d)} Xà phòng hóa hoàn toàn 8,06 gam tripalmitin bằng dung dịch NaOH vừa đủ, cô cạn thu được tối đa 8,44 gam muối sodium palmitate khan.\\[2pt]
\vspace{4pt}
\textbf{Câu 2.} Cho các thông tin về cấu tạo và tính chất của tinh bột và cellulose:\\[2pt]

\quad \textbf{a)} Tinh bột là polymer thiên nhiên gồm amylose (mạch không phân nhánh, xoắn lò xo) và amylopectin (mạch phân nhánh).\\[2pt]
\quad \textbf{b)} Hồ tinh bột tác dụng với dung dịch iodine ở nhiệt độ thường cho hợp chất màu xanh tím; khi đun nóng màu xanh tím biến mất, để nguội màu xanh tím xuất hiện trở lại.\\[2pt]
\quad \textbf{c)} Cellulose có 3 nhóm -OH tự do, khi phản ứng với nitric acid đặc có mặt sulfuric acid đặc xúc tác, đun nóng tạo thành cellulose trinitrate dùng làm thuốc súng không khói.\\[2pt]
\quad \textbf{d)} Cellulose là nguồn nguyên liệu chính trong công nghiệp sản xuất giấy, tơ sợi nhân tạo (tơ visco, tơ acetate) và vật liệu xây dựng.\\[2pt]

\vspace{4pt}
\noindent\textbf{\large\color{mainGreen}PHẦN III. Câu trắc nghiệm yêu cầu trả lời ngắn (2,0 điểm)}\\
\textit{Thí sinh trả lời từ Câu 1 đến Câu 4. Thí sinh ghi đáp số hoặc kết quả vào phiếu trả lời.}
\vspace{4pt}

\textbf{Câu 1.} Để tráng một số lượng gương soi có diện tích bề mặt $0,35\text{ m}^2$ với độ dày lớp bạc là $0,1\ \mu\text{m}$, người ta đun nóng dung dịch chứa $26,82\text{ gam}$ glucose với lượng dư thuốc thử Tollens. Biết khối lượng riêng của bạc là $10,49\text{ g/cm}^3$ và hiệu suất phản ứng tráng bạc đạt $80\%$ (tính theo glucose). Tính số lượng gương soi tối đa sản xuất được (làm tròn đến phần nguyên).\\[2pt]
\textbf{Trả lời:} \dotfill\\[4pt]
\vspace{2pt}
\textbf{Câu 2.} Hợp chất hữu cơ X là amine bậc một, no, đơn chức, mạch hở. Cho $9,0\text{ gam}$ X tác dụng với dung dịch HCl dư, cô cạn cẩn thận dung dịch sau phản ứng thu được $16,3\text{ gam}$ muối khan. Xác định số nguyên tử hydrogen trong phân tử amine X.\\[2pt]
\textbf{Trả lời:} \dotfill\\[4pt]
\vspace{2pt}
\textbf{Câu 3.} Cho m gam glycine tác dụng vừa đủ với $150\text{ mL}$ dung dịch HCl $1,0\text{ M}$ thu được dung dịch X. Cô cạn cẩn thận dung dịch X thu được lượng muối khan là bao nhiêu gam? (Kết quả làm tròn đến một chữ số thập phân).\\[2pt]
\textbf{Trả lời:} \dotfill\\[4pt]
\vspace{2pt}
\textbf{Câu 4.} Cho các dung dịch sau: glucose, saccharose, lòng trắng trứng (albumin), glycerol, ethylamine, Gly-Ala. Có bao nhiêu dung dịch phản ứng được với \ce{Cu(OH)2} ở điều kiện thường tạo thành dung dịch có màu xanh lam hoặc màu tím đặc trưng?\\[2pt]
\textbf{Trả lời:} \dotfill\\[4pt]
\vspace{2pt}

\vspace{4pt}
\noindent\textbf{\large\color{mainGreen}PHẦN IV. Tự luận (3,0 điểm)}\\
\textit{Thí sinh trình bày chi tiết lời giải các câu hỏi sau:}
\vspace{4pt}

Câu 1 (1,0 điểm). Một cơ sở sản xuất thủ công sử dụng 120 kg muối sodium của acid béo (thành phần chính của xà phòng) để sản xuất các bánh xà phòng thơm. Mỗi bánh xà phòng thành phẩm có khối lượng tịnh 92 gam và chứa 80\% khối lượng muối của acid béo (phần còn lại là các chất phụ gia, chất diệt khuẩn và hương liệu).\\
a) Viết phương trình hóa học của phản ứng điều chế xà phòng từ tripalmitin và dung dịch NaOH.\\
b) Tính số bánh xà phòng tối đa mà cơ sở trên sản xuất được (lấy kết quả nguyên).\\[2pt]
\par\noindent\dotfill\\
\noindent\dotfill\\
\noindent\dotfill\\
\noindent\dotfill\\
\noindent\dotfill\\[4pt]
\vspace{2pt}
Câu 2 (1,0 điểm). Một nhà máy chế biến đường mía ép 10 tấn mía tươi (chứa 13\% saccharose về khối lượng) thu được nước mía. Sau quá trình tinh chế với hiệu suất toàn bộ đạt 85\%, nhà máy thu được m kg đường saccharose tinh khiết.\\
a) Nêu đặc điểm cấu tạo của phân tử saccharose (gốc monosaccharide tạo thành, liên kết glycoside, dạng mạch tồn tại).\\
b) Xác định giá trị của m.\\[2pt]
\par\noindent\dotfill\\
\noindent\dotfill\\
\noindent\dotfill\\
\noindent\dotfill\\
\noindent\dotfill\\[4pt]
\vspace{2pt}
Câu 3 (1,0 điểm). Thủy phân hoàn toàn 14,6 gam dipeptide Gly-Ala bằng một lượng vừa đủ dung dịch NaOH, đun nóng nhẹ.\\
a) Viết phương trình hóa học của phản ứng thủy phân Gly-Ala trong dung dịch NaOH.\\
b) Cô cạn dung dịch sau phản ứng thu được m gam chất rắn khan. Tính giá trị của m.\\[2pt]
\par\noindent\dotfill\\
\noindent\dotfill\\
\noindent\dotfill\\
\noindent\dotfill\\
\noindent\dotfill\\[4pt]
\vspace{2pt}

\begin{center}
\textbf{---------- HẾT ----------}
\end{center}
\end{document}
